\documentclass[10pt,a4paper]{article}

% ============================================================
% B4 - POLITICA DE MANEJO DE DATOS DE GEOLOCALIZACION
% Documento autocontenido y compatible con pdfLaTeX en Overleaf.
% ============================================================

\usepackage[utf8]{inputenc}
\usepackage[T1]{fontenc}
\usepackage[spanish,es-nodecimaldot]{babel}
\usepackage{tgheros}
\renewcommand{\familydefault}{\sfdefault}
\usepackage[a4paper,top=1.85cm,bottom=1.8cm,left=1.9cm,right=1.9cm,headheight=20pt]{geometry}
\usepackage{xcolor}
\usepackage{tabularx,array,colortbl}
\usepackage{enumitem}
\usepackage[most]{tcolorbox}
\usepackage{fancyhdr}
\usepackage{lastpage}
\usepackage{microtype}
\usepackage{setspace}
\usepackage{titlesec}
\usepackage{hyperref}
\usepackage{xurl}
\usepackage{tikz}
\usepackage{ragged2e}

% ---------- PALETA ROUTTECH ----------
\definecolor{RTForest}{HTML}{25463B}
\definecolor{RTTerracotta}{HTML}{B45F45}
\definecolor{RTGold}{HTML}{D5A64A}
\definecolor{RTCream}{HTML}{F7F2E8}
\definecolor{RTLight}{HTML}{EEF3F0}
\definecolor{RTGray}{HTML}{4D5652}

% ---------- DATOS ----------
\newcommand{\Proyecto}{RoutTech: Sistema de Movilidad Colaborativa para la Comunidad Universitaria de la UTEQ}
\newcommand{\Asignatura}{ISR-401 -- Ingeniería de Requisitos}
\newcommand{\Periodo}{2026--2027 PPA}
\newcommand{\CategoriaEtica}{Categoría B -- Datos personales}
\newcommand{\Repositorio}{https://github.com/AnderJM3/RoutTech_ISR401}
\newcommand{\FechaDocumento}{29 de julio de 2026}
\newcommand{\VersionDocumento}{2.0}
\newcommand{\Paralelo}{Cuarto nivel -- Software B}
\newcommand{\Mayra}{CORDOVA CARREÑO MAYRA LUCILA}
\newcommand{\Anderson}{NARANJO FLORES ANDERSON JEAMPIERE}
\newcommand{\Docente}{Ing. Gleiston Guerrero Ulloa, PhD}

% ---------- TABLAS ----------
\newcolumntype{Y}{>{\RaggedRight\arraybackslash}X}
\newcolumntype{C}[1]{>{\centering\arraybackslash}p{#1}}
\newcolumntype{L}[1]{>{\RaggedRight\arraybackslash}p{#1}}
\renewcommand{\arraystretch}{1.20}
\setlength{\tabcolsep}{4.5pt}

% ---------- ENCABEZADO Y PIE ----------
\pagestyle{fancy}
\fancyhf{}
\fancyhead[L]{\small\textbf{UTEQ | FCC | ISR-401}}
\fancyhead[R]{\small\textbf{B4 -- Política de geolocalización}}
\fancyfoot[L]{\small Paquete ético RoutTech}
\fancyfoot[R]{\small Página \thepage\ de \pageref{LastPage}}
\renewcommand{\headrulewidth}{0.8pt}
\renewcommand{\footrulewidth}{0.4pt}
\renewcommand{\headrule}{\hbox to\headwidth{\color{RTTerracotta}\leaders\hrule height \headrulewidth\hfill}}
\renewcommand{\footrule}{\hbox to\headwidth{\color{RTGold}\leaders\hrule height \footrulewidth\hfill}}

% ---------- TITULOS ----------
\titleformat{\section}
  {\large\bfseries\color{RTForest}}
  {\thesection.}{0.45em}{}
  [\vspace{1pt}\color{RTGold}\titlerule]
\titlespacing*{\section}{0pt}{0.65em}{0.30em}

\hypersetup{
  colorlinks=true,
  linkcolor=RTForest,
  urlcolor=RTTerracotta,
  pdftitle={B4 Politica de Geolocalizacion RoutTech},
  pdfauthor={Equipo RoutTech}
}
\urlstyle{same}
\setlength{\parindent}{0pt}
\setlength{\parskip}{2.5pt}
\setstretch{1.06}
\setlist[enumerate]{leftmargin=1.55em,itemsep=2pt,topsep=2pt}

\begin{document}
\justifying

% ======================= PORTADA =======================
\thispagestyle{empty}
\begin{tikzpicture}[remember picture,overlay]
  \fill[RTForest] (current page.north west) rectangle ([yshift=-4.6cm]current page.north east);
  \fill[RTTerracotta] ([yshift=-4.6cm]current page.north west) rectangle ([yshift=-4.95cm]current page.north east);
\end{tikzpicture}

\vspace*{0.65cm}
\begin{center}
{\color{white}\Large\bfseries UNIVERSIDAD TÉCNICA ESTATAL DE QUEVEDO}

\vspace{1.75cm}
{\large\bfseries\color{RTForest} Facultad de Ciencias de la Computación\\[4pt]}
{\normalsize\color{RTGray} Carrera de Ingeniería de Software}

\vspace{0.8cm}

{\Huge\bfseries\color{RTForest} POLÍTICA DE GEOLOCALIZACIÓN\\[8pt]}
{\Large\bfseries\color{RTTerracotta} B4 -- Categoría B}

\vspace{1.0cm}

\begin{tcolorbox}[width=0.90\textwidth,colback=RTCream,colframe=RTGold,boxrule=1pt,arc=2mm]
\centering
{\large\bfseries \Proyecto}\\[7pt]
\textbf{Clasificación ética:} \CategoriaEtica\\
\textbf{Asignatura:} \Asignatura\\
\textbf{Periodo académico:} \Periodo
\end{tcolorbox}

\vfill
\begin{tabular}{rl}
\textbf{Documento:} & B4 -- Política de manejo de datos de geolocalización\\
\textbf{Versión:} & \VersionDocumento\\
\textbf{Estado:} & Versión 2.0 para firma\\
\textbf{Paralelo:} & \Paralelo\\
\textbf{Repositorio:} & \href{\Repositorio}{RoutTech\_ISR401}\\
\textbf{Fecha:} & \FechaDocumento
\end{tabular}

\vspace{0.75cm}
{\small Quevedo -- Los Ríos -- Ecuador}
\end{center}
\clearpage

% ======================= CONTENIDO =======================
\section{Alcance}

Esta política regula el tratamiento de datos de ubicación en las actividades de investigación, validación, pruebas, demostraciones y publicación del proyecto RoutTech.

\section{Reglas de manejo}

\begin{enumerate}
  \item Las ubicaciones se tratarán de forma agregada o generalizada y no se asociarán públicamente con una persona identificada o identificable.
  \item Los puntos de origen y destino se expresarán como campus, sector, barrio o parroquia; no se registrarán ni publicarán direcciones domiciliarias exactas.
  \item No se publicarán coordenadas precisas, mapas con trayectorias individuales ni historiales de desplazamiento de conductores o pasajeros.
  \item Las rutas utilizadas en el MVP, pruebas, capturas y demostraciones serán sintéticas o estarán previamente disociadas.
  \item Durante el desarrollo académico, entrevistas y sesiones de validación no se activará seguimiento de ubicación en tiempo real.
  \item Una futura función de ubicación requerirá consentimiento explícito, temporal y revocable del usuario, y solo se mantendrá activa durante la operación necesaria.
  \item La negativa o retiro del permiso de ubicación no impedirá el acceso a las funciones que no dependan de este dato.
  \item Los datos de ubicación no se utilizarán para vigilancia académica, disciplinaria, laboral, comercial ni para elaborar perfiles individuales.
  \item Los reportes y paneles mostrarán únicamente resultados agregados y omitirán grupos pequeños que puedan facilitar la reidentificación.
  \item Las evidencias visuales y los archivos publicados excluirán domicilios, placas, coordenadas, rutas personales y metadatos de ubicación.
\end{enumerate}

\section{Controles de verificación}

\rowcolors{2}{RTLight}{white}
\noindent\begin{tabularx}{\textwidth}{L{5.0cm} Y}
\rowcolor{RTForest}
\textcolor{white}{\textbf{Control}} & \textcolor{white}{\textbf{Evidencia esperada}}\\
Datos sintéticos en el MVP & Casos de prueba, base de demostración y capturas con rutas ficticias.\\
Permiso explícito y revocable & Requisito, escenario Gherkin y prueba funcional del consentimiento y retiro.\\
Generalización de ubicaciones & Revisión de campos, mapas y archivos antes de su uso o publicación.\\
Ausencia de rutas individuales & Auditoría del repositorio y de los documentos de divulgación.\\
Desactivación del seguimiento & Configuración y pruebas que demuestren la ausencia de geolocalización en tiempo real.\\
\end{tabularx}

\section{Firmas}

\vspace{0.25cm}
\noindent\begin{tabularx}{\textwidth}{Y C{4.1cm} C{2.9cm}}
\rowcolor{RTTerracotta}
\textcolor{white}{\textbf{Nombre}} &
\textcolor{white}{\textbf{Firma}} &
\textcolor{white}{\textbf{Fecha}}\\
\Mayra & \rule{3.7cm}{0.4pt} & \rule{2.4cm}{0.4pt}\\[4pt]
\Anderson & \rule{3.7cm}{0.4pt} & \rule{2.4cm}{0.4pt}\\[4pt]
\Docente & \rule{3.7cm}{0.4pt} & \rule{2.4cm}{0.4pt}\\
\end{tabularx}

\end{document}
